\subsection{Bisection and complete enumeration}

Budgets are listed from largest to smallest, so the desired answer is
the last acceptable grid index. Put $J=|A|$. After computing the sample
reference once, the implementation searches as follows:
\par\noindent\begin{minipage}{\linewidth}
\begin{enumerate}[leftmargin=*,itemsep=0.15em]
\item Set $l=0$, $h=J-1$. Zero is the fallback.
\item While $l<h$, test $j=\lfloor(l+h+1)/2\rfloor$ by running
$A[j]$ and computing $D_S(A[j])$.
\item If $D_S(A[j])\leq\epsilon$, set $l=j$; otherwise set $h=j-1$.
\item Return $A[l]$. Apply it before reuse.
\end{enumerate}
\end{minipage}\par\medskip
The eight-point grid needs at most three trial searches plus the
reference. Enumeration instead tries the reference and all seven other
budgets on the same sample. It finds the minimum even if accepted
budgets are separated by rejected ones; bisection is guaranteed to find it when
all accepted indices form a prefix. Both timers include the reference
search (\cref{sec:a-budget-contract}).

\subsection{A sufficient condition for monotonicity}

A sufficient structure is nested shortlists: every candidate retained at
a smaller budget is also present at a larger one. With fixed refined
scores and final tie ordering, a reference answer that survives the
shortlist must survive final selection. The proposition makes this
precise and shows why disagreement then falls as the budget grows.

\par\noindent\begin{minipage}{\linewidth}
\begin{proposition}[Nested survivors preserve reference membership]
\label{prop:6-monotonicity}
For one query, assume fixed routing and survivor sets
$C_\alpha\subseteq C_\beta\subseteq C_{\alpha_{\rm ref}}$ whenever
$\alpha\leq\beta\leq\alpha_{\rm ref}$. Assume at least $k$ valid
survivors, fixed refined scores and one deterministic total order for
final ties. Each result is the exact top-$k$ in that order. Then
\begin{equation}
 |\mathcal R_\alpha\cap\mathcal R_{\rm ref}|
 =|C_\alpha\cap\mathcal R_{\rm ref}|,
 \qquad D_S(\alpha)\geq D_S(\beta).
 \label{eq:6-monotonicity}
\end{equation}
The latter statement holds when these conditions hold for every sampled
query.
\end{proposition}
\begin{proof}
For $x\in\mathcal R_{\rm ref}\cap C_\alpha$, the common total order gives
\begin{equation*}
 |\{y\in C_\alpha:y\prec x\}|
 \leq|\{y\in C_{\alpha_{\rm ref}}:y\prec x\}|<k.
\end{equation*}
Hence $x\in\mathcal R_\alpha$. Together with
$\mathcal R_\alpha\subseteq C_\alpha$, this proves the intersection
identity in \cref{eq:6-monotonicity}. Nestedness gives
$|C_\alpha\cap\mathcal R_{\rm ref}|
 \leq|C_\beta\cap\mathcal R_{\rm ref}|$.
Summing missing reference counts over queries proves
$D_S(\alpha)\geq D_S(\beta)$.
\end{proof}
\end{minipage}\par\medskip
Exact primary top-$c$ under one fixed key gives nested prefixes. This
makes the reference-agreement test monotone under the proposition's
other premises; GT recall and runtime need not be monotone.

\par\medskip\noindent\begin{minipage}{\linewidth}
\captionsetup{hypcap=false}
\centering
\captionof{table}[Nested and non-nested candidate pools]{Nested and non-nested candidate pools.}
\label{tab:6-nesting}
\begingroup
\small\setlength{\tabcolsep}{6pt}
\begin{tabular}{@{}lcccc@{}}
\toprule
Case & $C_1$ & $C_2$ & $C_3$ & $|C_i\cap\mathcal R_{\rm ref}|$, $i=1,2,3$ \\
\midrule
Nested & $\{b\}$ & $\{b,c\}$ & $\{a,b,c\}$ & $0,\ 0,\ 1$ \\
Non-nested & $\{a\}$ & $\{b,c\}$ & $\{a,b,c\}$ & $1,\ 0,\ 1$ \\
\bottomrule
\end{tabular}
\endgroup

\end{minipage}\par\medskip
\Cref{tab:6-nesting} gives a $k=1$ example with refined order
$a\prec b\prec c$ and reference $\{a\}$. Nested pools preserve reference
membership. If primary scores all tie, changing cutoff choices can
instead make the middle pool exclude an answer retained by the smaller
pool; overlap falls from 1 to 0 before returning to 1.

\subsection{A corollary for the evaluated grouped path}

Within one calibration sample, $S$, routing and the scan thread count
are fixed. The scheduler chooses group count without using $c$, so
reference and trial budgets use the same scan path and LUT sums.
Fixed $B_r$, sixteen-lane refinement and final
$(\refineddist,\mathrm{ID})$ ordering keep each candidate's refined
score and tie order independent of the budget. Grouped selection then
supplies the nested primary shortlists required above.

\par\addvspace{7pt}\noindent\begin{minipage}{\linewidth}
\begin{corollary}[Grouped calibration]
\label[corollary]{cor:6-grouped-calibration}
For the evaluated grouped configuration, under the representation and
shared-bound contracts of Chapters 4 and 5, primary survivors are nested
across budgets. Consequently \Cref{prop:6-monotonicity} applies and
bisection returns the same minimum accepted grid budget as enumeration.
\end{corollary}
\begin{proof}
\Cref{lem:5-buffer,prop:5-local-coverage,prop:5-shared-bound} give exact
primary top-$c$ under the common key $(\primarydist,\mathrm{position})$.
Its prefixes nest as $c$ grows. Fixed refinement/final order supplies \cref{prop:6-monotonicity}'s remaining premises.
\end{proof}
\end{minipage}\par\medskip
The corollary covers 2,176 calibration samples. Another 128 samples,
OpenAI-3072 at $S=128$, use distance-only selection, where cutoff ties
can break nesting. In the recorded experiment, all 2,304 curves were
monotone and bisection matched enumeration
(\cref{sec:7-budget-trace}). This validates the recorded choices, while
the corollary supplies the guarantee for the grouped path.

\paragraph{Full-pool coverage.}
Two full-pool settings use grouped scans on Vogue; the other ten use
$\gscan=0$. All twelve match enumeration. The ten distance-only settings
still lie outside \cref{cor:6-grouped-calibration}: they need a common
cutoff order, or no tied reference answers that cross the cutoff.
Using the same path for calibration and testing does not establish nesting.

\paragraph{A rank-based extension.}
Under the same nesting and fixed-score premises, storing the primary
rank of each reference answer would reveal when each budget loses it:
\begin{equation}
 D_S(\alpha)=\sum_{q\in\mathcal Q_S}\sum_{x\in\mathcal R_{\rm ref}(q)}
 \mathbf1\{\operatorname{rank}_{\rm P}(x)>c(\alpha)\}.
 \label{eq:6-rank-curve}
\end{equation}
This could produce the whole agreement curve from exported reference
ranks instead of repeated searches; it is not implemented.
